\documentclass[runningheads]{llncs}

\usepackage[T1]{fontenc}
\usepackage{amsmath,amssymb,amsfonts}
\usepackage{booktabs}
\usepackage{multirow}
\usepackage{xcolor}
\usepackage{hyperref}
\usepackage{algorithm}
\usepackage{algpseudocode}

\begin{document}

\title{Privacy-Preserving Decision Tree Inference under CKKS:\\A Precision Impact Analysis}

\titlerunning{CKKS Decision Tree Inference: Precision Impact}

\author{Madicke-Diadji Mbodj\inst{1}}

\authorrunning{M.-D. Mbodj}

\institute{[Institution Name]\\
\email{diazlinkk@gmail.com}}

\maketitle

%% =========================================================
\begin{abstract}
Privacy-preserving machine learning inference via Fully Homomorphic
Encryption (FHE) allows a server to classify encrypted client data without
ever accessing plaintext.  Decision trees are attractive inference models
for this setting owing to their simplicity and interpretability, yet their
discontinuous branching structure is fundamentally at odds with the smooth
arithmetic required by modern FHE schemes.  We present a complete
implementation of private decision tree inference combining the HBDT-SumPath
constant-depth algorithm with an \emph{adaptive} polynomial approximation of
the step function, where each internal tree node is assigned a polynomial
degree proportional to the difficulty of its splitting decision -- quantified
by the \emph{minimum feature gap} observed on training data.  Our central
contribution is a systematic analysis of the two independent sources of
accuracy loss that arise in this pipeline: (i) approximation error due to
smooth replacement of the step function, and (ii) residual noise introduced
by the CKKS homomorphic encryption scheme.  We evaluate on seven benchmark
datasets, quantify the relationship between minimum gap and observed accuracy
degradation, and identify when each source dominates.  On several datasets,
the final adaptive HE pipeline matches the adaptive cleartext baseline
exactly; on harder cases, a substantial part of the loss is due to
\emph{unnormalised comparison logits} rather than CKKS itself.  After
introducing node-wise logit normalisation, the adaptive pipeline recovers
the hard-tree accuracy on Breast, Heart, and Steel in both cleartext and
CKKS.  In the latest runs, the main residual CKKS-only loss remains on
Cancer, where the small feature gap still induces a $3.1$-point drop
between adaptive cleartext and adaptive HE.  Relative to the historical
Akavia-style SEAL baseline visible in our local benchmark plots, our
implementation is also markedly faster on the overlapping standard
datasets.  All experiments use Microsoft SEAL v4.1.2 at the 128-bit
security level.
\end{abstract}

\keywords{Homomorphic Encryption \and CKKS \and Decision Tree \and
Privacy-Preserving Inference \and Polynomial Approximation \and
Data Privacy}

%% =========================================================
\section{Introduction}
\label{sec:intro}

The widespread deployment of machine learning models as a service raises a
fundamental tension between utility and privacy.  A client wishing to obtain
a classification result must either reveal its sensitive input to the server,
or perform inference locally -- requiring possession of the model.
\emph{Private inference} dissolves this dilemma: the client encrypts its
data, the server evaluates the model directly on ciphertext, and the client
decrypts the result.  Neither party learns the other's secret.

Decision trees occupy a privileged position in privacy-preserving machine
learning (PPML) owing to their compactness, interpretability, and widespread
deployment in regulated domains such as healthcare, credit scoring, and
intrusion detection~\cite{bos2014private,tai2017mining}.  A depth-$D$ tree
can be represented in $\mathcal{O}(2^D)$ internal nodes, each performing a
scalar comparison $x[\mathit{feat}] \leq \theta$.  Under FHE, however,
comparison is not a native operation: it must be approximated by a smooth
polynomial, since every FHE scheme over the integers or reals supports only
addition and multiplication.

Among current FHE schemes, \textbf{CKKS}~\cite{cheon2017homomorphic}
(Cheon-Kim-Kim-Song) is the scheme of choice for real-valued inference.
CKKS encodes vectors of floating-point numbers in ciphertext slots and
supports batched SIMD arithmetic, but introduces a controlled \emph{noise}
after each multiplication.  This noise is bounded and manageable -- yet when
the polynomial approximation of the comparison function evaluates near its
transition region (i.e., when the true input is close to the threshold
$\theta$), even a small CKKS noise perturbation can flip the effective
branch decision and thus degrade classification accuracy.

\textbf{Contributions.}  This paper makes the following contributions:
\begin{enumerate}
  \item We implement a complete private decision tree inference system
        combining the HBDT-SumPath depth-reduction
        algorithm~\cite{shin2022privacy} with per-node adaptive polynomial
        approximation, using Microsoft SEAL v4.1.2.
  \item We introduce and validate the \emph{minimum feature gap}
        (\textit{min\_gap}) as an interpretable metric that predicts the
        sensitivity of a given tree and dataset to both sources of accuracy
        loss.
  \item We show that for several pre-trained benchmark trees, local
        \emph{node-wise logit normalisation} is necessary for the adaptive
        approximation to reflect the true decision difficulty of each node.
  \item We provide the \textbf{first systematic disentanglement} of
        approximation error (visible in cleartext soft evaluation) from CKKS
        noise error (the additional degradation in HE evaluation), across
        seven benchmark datasets.
  \item We characterise three precision regimes as a function of
        \textit{min\_gap} and provide concrete parameter recommendations to
        mitigate accuracy loss in each regime.
\end{enumerate}

The remainder of the paper is organised as follows.
Section~\ref{sec:background} reviews CKKS, decision trees, and the
SumPath algorithm.  Section~\ref{sec:system} describes the system
architecture.  Section~\ref{sec:precision} develops the precision analysis.
Section~\ref{sec:experiments} reports experimental results.
Section~\ref{sec:related} discusses related work, and
Section~\ref{sec:conclusion} concludes.

%% =========================================================
\section{Background}
\label{sec:background}

\subsection{CKKS Homomorphic Encryption}

CKKS~\cite{cheon2017homomorphic} is a levelled FHE scheme over the ring
$\mathbb{Z}[X]/(X^N+1)$ that encodes a vector of $N/2$ complex (or real)
values into a single polynomial ciphertext.  Arithmetic over ciphertext
closely mimics floating-point arithmetic: additions are exact and
multiplications introduce a rounding error of magnitude
$\varepsilon_{\mathrm{mult}} \approx \Delta^{-1}$ per level, where $\Delta$
is the \emph{scale} factor chosen at key generation.

In Microsoft SEAL~\cite{chen2020simple}, the scheme is parameterised by:
\begin{itemize}
  \item $N$ -- the polynomial modulus degree (power of 2); we use
        $N = 2^{14} = 16{,}384$, giving $N/2 = 8{,}192$ SIMD slots;
  \item $\Delta = 2^{40}$ -- the encoding scale, providing approximately
        12 decimal digits of precision per ciphertext level;
  \item $L$ -- the multiplicative depth budget (number of available levels
        before the ciphertext noise floor becomes too large for correct
        decryption); we configure $L = 12$.
\end{itemize}

After a sequence of $\ell$ multiplications the accumulated noise satisfies:
\begin{equation}
  \varepsilon_{\ell} \;\approx\; \ell \cdot N^{1/2} \cdot \Delta^{-1}
  \;\approx\; \ell \cdot 10^{-7},
  \label{eq:ckks_noise}
\end{equation}
with $N = 2^{14}$ and $\Delta = 2^{40}$.  Concretely, for a depth-6
evaluation $\ell = 6$, Eq.~\eqref{eq:ckks_noise} gives
$\varepsilon_6 \approx 6 \times 10^{-7}$, well below a gap threshold of
$\delta = 0.05$ in typical settings.  The bound tightens, however, when
polynomial evaluation itself occupies most of the depth budget and the
residual noise after evaluation approaches the minimum feature gap.

\subsection{Decision Tree Inference}

A binary decision tree of depth $D$ over $d$ features partitions the input
space $\mathbb{R}^d$ through a sequence of threshold comparisons.  Each
internal node $N_i$ tests $x[\mathit{feat}_i] \leq \theta_i$; every leaf
stores a class label.  A sample traverses the tree from root to leaf by
following the satisfied branch at each node, and is classified with the
label stored at the reached leaf.

\textbf{The comparison problem.}  Computing $\mathbf{1}[x \leq \theta]$
over a CKKS ciphertext requires replacing the non-differentiable
indicator $\mathbf{1}[t \geq 0]$ (with $t = \theta - x$) by a polynomial
$\varphi(t)$ satisfying:
\begin{equation}
  \varphi(t) \approx
  \begin{cases} 1 & t > \delta \\ 0 & t < -\delta \end{cases},
  \qquad \varphi(t) \in [0,1] \;\forall t.
  \label{eq:step_approx}
\end{equation}
The interval $(-\delta, \delta)$ is a \emph{dead window} within which the
approximation is not required to be accurate.  Any input sample with
$|\theta_i - x[\mathit{feat}_i]| < \delta$ lies within the dead window and
may be misclassified independently of the polynomial degree.

\subsection{HBDT-SumPath Algorithm}
\label{sec:sumpath}

A naive evaluation of all $2^D - 1$ node comparisons in an FHE circuit
consumes $\mathcal{O}(D \cdot \deg \varphi)$ multiplicative levels.  For
$D = 4$ and $\deg \varphi = 8$ this already requires 32 levels -- near the
practical limit of $L = 12$ for $N = 2^{14}$.

Shin \emph{et al.}~\cite{shin2022privacy} proposed HBDT-SumPath, which
reformulates tree evaluation as follows.  Define
$I_i = 1 - \varphi(\theta_i - x[\mathit{feat}_i])$; then $I_i \approx 0$
means the left branch is taken at node $N_i$, and $I_i \approx 1$ means the
right branch.  For a leaf $f_k$ reached by the path
$\{N_1, \ldots, N_D\}$, let:
\begin{equation}
  S(f_k) \;=\; \sum_{i : N_i \in \mathit{path}(f_k)} I_i.
  \label{eq:sumpath}
\end{equation}
The correct leaf has $S = 0$ (all comparisons satisfied); any other leaf has
$S \geq 1$.  The SumPath reduction converts multiplicative depth from
$\mathcal{O}(D)$ to $\mathcal{O}(1)$: once the per-node indicators are
evaluated, the sum over a path is a \emph{linear} operation (additions and
slot rotations) consuming no additional multiplicative levels.

%% =========================================================
\section{System Architecture}
\label{sec:system}

The proof-of-concept follows a simple five-stage pipeline: training or
loading a tree, cleartext soft evaluation for calibration, client-side
encryption, server-side HE evaluation under CKKS, and client-side
decryption of the predicted label.  Along this path, the implementation
also exports tree statistics such as \textit{min\_gap}, derives the
per-node degree map, and records ciphertext-side scores for evaluation.

\subsection{Polynomial Approximation of the Step Function}
\label{sec:softapprox}

We approximate $\mathbf{1}[t \geq 0]$ on $[-2, 2]$ by a degree-$n$ odd
polynomial $p_n(t)$ computed by weighted least-squares:
\begin{equation}
  p_n \;=\; \arg\min_{p \in \mathcal{P}_n}
    \int_{-2}^{2} \bigl(\mathbf{1}[t \geq 0] - p(t)\bigr)^2 \, w(t)\,dt,
  \label{eq:wls}
\end{equation}
where $w(t) = 0$ for $|t| < \delta$ (dead window, $\delta = 0.05$) and
$w(t) = 1$ elsewhere.  The integral is discretised over 2{,}000 uniform
points and solved by Gaussian elimination.  The final approximation is
$\varphi(t) = \mathtt{clamp}(p_n(t),\, 0,\, 1)$.

\textbf{Evaluation under CKKS.}  A degree-$n$ polynomial evaluated naively
requires $n-1$ multiplications, consuming $n-1$ levels.  We use the
Baby-Step Giant-Step (BSGS) algorithm~\cite{han2020better}, which reduces
this to $\lceil \log_2 n \rceil + \lceil n/k \rceil - 2$ levels for optimal
baby-step size $k = \lfloor \sqrt{n} \rfloor$.  For our degree set
$\{4, 8, 16, 32\}$ this yields depths $\{2, 3, 4, 5\}$ respectively,
all within the budget $L = 12$ even after SumPath overhead.

\subsection{Adaptive Degree Selection via Minimum Feature Gap}
\label{sec:adaptive}

The core insight motivating our adaptive strategy is that different
internal nodes of the same tree exhibit very different decision difficulty:
a node that always sees training samples far from its threshold needs only a
coarse polynomial, while a node whose training samples cluster near the
threshold requires a fine one.

\textbf{Definition (minimum feature gap).}  For internal node $N_i$ with
threshold $\theta_i$ on feature $\mathit{feat}_i$:
\begin{equation}
  \mathit{min\_gap}(N_i) \;=\; \min_{x \in \mathcal{X}\,:\,N_i \text{ active}}
    \bigl|x[\mathit{feat}_i] - \theta_i\bigr|,
  \label{eq:mingap}
\end{equation}
where ``$N_i$ active'' means the training sample $x$ reaches $N_i$.
We collect $\mathit{min\_gap}$ for every node via scikit-learn's
\texttt{decision\_path()} API during training and export the values to a
configuration file consumed by the C++ inference engine.

\textbf{Degree assignment rule.}  Let $g = \mathit{min\_gap}(N_i)$:
\begin{align}
  n(g) = \begin{cases}
    4  & g \geq 0.08   \\
    8  & 0.04 \leq g < 0.08 \\
    16 & 0.015 \leq g < 0.04 \\
    32 & g < 0.015 \text{ or } g = 0
  \end{cases}
  \label{eq:degmap}
\end{align}
This rule ensures that nodes whose minimum gap falls within the dead window
$\delta = 0.05$ are assigned the highest degree, maximising approximation
quality where it matters most.

\textbf{Node-wise logit normalisation.}  For the pre-trained SortingHat
trees, we found that the raw comparison logits $\theta_i - x[\mathit{feat}_i]$
can have very different scales across nodes, which makes the adaptive degree
heuristic unstable when it is driven by an unnormalised gap alone.  We
therefore calibrate each internal node with
\begin{equation}
  \mathit{norm}_i = \max_{x \in \mathcal{X}_{\mathrm{calib}}}
  \bigl|x[\mathit{feat}_i] - \theta_i\bigr|,
\end{equation}
and evaluate the soft comparator on the normalised logit
$(\theta_i - x[\mathit{feat}_i]) / \mathit{norm}_i$.  The adaptive degree is
then selected from the normalised gap
$\mathit{min\_gap}(N_i) / \mathit{norm}_i$.  This calibration step is what
allows the adaptive pipeline to recover the hard-tree baseline on Breast,
Heart, and Steel.

\subsection{CKKS Parameters and Vertical Packing}

The implementation uses Microsoft SEAL 4.1.2 and configures CKKS at runtime.
The polynomial modulus degree $N$ is selected automatically from
$\{2^{13}, 2^{14}, 2^{15}\}$, i.e., $\{8{,}192, 16{,}384, 32{,}768\}$,
according to the slot budget and the coefficient modulus chain.  The latter
is instantiated as $\{60, \underbrace{s,\ldots,s}_{L\ \mathrm{times}}, 60\}$,
where $s \in \{40,45,50\}$ is the scale modulus size of the active profile
and $L \geq 10$ is the target multiplicative depth.  The effective CKKS
scale is $\Delta = 2^s$.  For the larger datasets used in our main runs,
this typically yields $N = 2^{14}$ and $s = 40$, consistent with a 128-bit
security configuration under the HE standard~\cite{albrecht2021homomorphic}.

\textbf{Vertical packing (SIMD batching).}  Each input sample of $d$ features
is encoded over $n_{\max}$ buckets per feature, where $n_{\max}$ is profile
dependent ($4$ for the small profile, $2$ otherwise), occupying
$d \cdot n_{\max}$ slots.  Multiple samples are packed into a single
ciphertext by vertical concatenation.  For example, with $N/2 = 8{,}192$
slots and $d = 30$ features under the large profile ($n_{\max}=2$), one
sample occupies $60$ slots.  The implementation further caps the effective
packed batch size to at most 32 samples in order to control memory usage and
the overhead of Galois rotations and key switching.

%% =========================================================
\section{Precision Impact Analysis}
\label{sec:precision}

This section develops the main analytical contribution of the paper:
a decomposition of the accuracy loss observed in private inference into two
independent and interpretable sources.

\subsection{Two Sources of Accuracy Loss}
\label{sec:two_sources}

Let $\mathcal{A}_{\mathrm{hard}}$ denote the accuracy of standard (hard)
cleartext inference, $\mathcal{A}_{\mathrm{soft}}$ the accuracy of cleartext
soft inference (using $\varphi$ instead of $\mathbf{1}$), and
$\mathcal{A}_{\mathrm{HE}}$ the accuracy of HE inference.  We define:
\begin{align}
  \Delta_{\mathrm{approx}} &\;=\; \mathcal{A}_{\mathrm{hard}}
                                  - \mathcal{A}_{\mathrm{soft}}, \label{eq:delta_approx} \\
  \Delta_{\mathrm{CKKS}}   &\;=\; \mathcal{A}_{\mathrm{soft}}
                                  - \mathcal{A}_{\mathrm{HE}}.   \label{eq:delta_ckks}
\end{align}

\textbf{Source 1: Approximation error} ($\Delta_{\mathrm{approx}}$).
This error arises whenever a sample falls within the dead window of some
active node: $|x[\mathit{feat}_i] - \theta_i| < \delta$.  In this region
$\varphi(\theta_i - x[\mathit{feat}_i]) \approx 0.5$, rendering the branch
decision stochastic.  This error is \emph{intrinsic to the approximation
strategy} and is fully visible in cleartext.  It can be reduced by
decreasing $\delta$ or increasing the polynomial degree $n$.

\textbf{Source 2: CKKS noise error} ($\Delta_{\mathrm{CKKS}}$).
This error arises from the residual noise $\varepsilon_\ell$
(Eq.~\eqref{eq:ckks_noise}) accumulated during homomorphic polynomial
evaluation.  For a sample with $|\theta_i - x| \gg \varepsilon_\ell$,
the noise is absorbed by the saturation of $\varphi$ and the branch
decision is unchanged.  The noise becomes decisive only when
$|\theta_i - x| \lesssim \varepsilon_\ell$, which occurs for samples with
$\mathit{min\_gap}$ near or below $\varepsilon_\ell \approx 10^{-6}$.

\begin{remark}
$\Delta_{\mathrm{approx}}$ and $\Delta_{\mathrm{CKKS}}$ are not additive in
general: a sample mis-classified by the approximation in cleartext is not
counted again in $\Delta_{\mathrm{CKKS}}$.  However, they are
\emph{statistically independent} for datasets where the two error populations
do not overlap.
\end{remark}

\subsection{Three Precision Regimes}
\label{sec:regimes}

Table~\ref{tab:regimes} summarises the three regimes as a function of
$\mathit{min\_gap}$ relative to the dead window $\delta$ and the CKKS noise
floor $\varepsilon$.

\begin{table}[t]
\centering
\caption{Precision regimes as a function of minimum feature gap.
$\delta = 0.05$, $\varepsilon \approx 10^{-6}$ (depth 6, $\Delta = 2^{40}$,
$N = 2^{14}$).}
\label{tab:regimes}
\begin{tabular}{@{}llll@{}}
\toprule
Regime & Condition & Dominant loss & Typical $\Delta_{\mathrm{CKKS}}$ \\
\midrule
I   (safe)    & $\mathit{min\_gap} \geq \delta$                  & None              & $< 1\%$     \\
II  (approx.) & $\varepsilon \leq \mathit{min\_gap} < \delta$    & $\Delta_{\mathrm{approx}}$ & $\approx 0\%$ in our runs \\
III (critical)& $\mathit{min\_gap} < 5 \times 10^{-3}$          & Approx.\ and/or CKKS & $0$--$22.7\%$ (unnormalised baseline) \\
\bottomrule
\end{tabular}
\end{table}

\textbf{Regime I.}  When $\mathit{min\_gap} \geq \delta$, all training
samples are outside the dead window.  The polynomial $\varphi$ saturates
to 0 or 1 before CKKS noise can accumulate, and both
$\Delta_{\mathrm{approx}}$ and $\Delta_{\mathrm{CKKS}}$ are negligible.

\textbf{Regime II.}  When $\mathit{min\_gap}$ falls between $\varepsilon$ and
$\delta$, some samples lie inside the dead window and experience
approximation error.  CKKS noise is still below the gap and does not add to
the cleartext loss.  This regime is correctly handled by increasing the
polynomial degree (reducing the effective dead window at the cost of
multiplicative depth).

\textbf{Regime III.}  When $\mathit{min\_gap} \ll \delta$ and in particular
$\mathit{min\_gap} \lesssim 5 \times 10^{-3}$, CKKS noise can become
comparable to the gap at critical nodes, producing additional HE-only errors
beyond those already visible in cleartext.  In this regime, accuracy
degradation can be substantial and requires either (a) increasing $\Delta$ to
$2^{50}$ to lower the noise floor, (b) increasing $N$ to $2^{15}$ to gain
additional depth budget, or (c) retraining the tree with a
$\mathit{min\_gap}$-regularisation objective.

\subsection{Noise Propagation Through SumPath}
\label{sec:noise_prop}

Under SumPath, the per-leaf score $S(f_k) = \sum_{i \in \mathit{path}} I_i$
accumulates noise linearly.  Let $\hat{I}_i = I_i + e_i$ be the noisy HE
indicator, with $|e_i| \leq \varepsilon_\ell$.  Then:
\begin{equation}
  \hat{S}(f_k) \;=\; S(f_k) + \sum_{i \in \mathit{path}} e_i, \qquad
  \bigl|\hat{S}(f_k) - S(f_k)\bigr| \;\leq\; D \cdot \varepsilon_\ell,
  \label{eq:noise_sum}
\end{equation}
where $D$ is the tree depth.  Since $S(f_k) \in \{0, 1, \ldots, D\}$ and
correct classification requires $S(f^*) = 0$ uniquely, a misclassification
occurs when $\hat{S}(f^*) \geq 1$ or $\hat{S}(f') \leq 0$ for some
$f' \neq f^*$.  With $D = 4$ and $\varepsilon_\ell \approx 10^{-6}$:
\begin{equation}
  D \cdot \varepsilon_\ell \;=\; 4 \times 10^{-6} \;\ll\; 1,
\end{equation}
meaning SumPath noise propagation is negligible compared to the
indicator-level approximation errors.  The dominant precision loss therefore
stems from individual node evaluations, not from path aggregation.

%% =========================================================
\section{Experimental Evaluation}
\label{sec:experiments}

\subsection{Setup}

\textbf{Datasets.}  We evaluate on seven benchmark datasets:
Iris ($d=4$, 3 classes), Wine ($d=13$, 3 classes), Cancer ($d=30$, 2
classes), Digits ($d=64$, 10 classes), Breast ($d=30$, 2 classes), Heart
($d=13$, 2 classes), and Steel ($d=33$, 2 classes).  For Iris, Wine, Cancer,
and Digits we train a scikit-learn \texttt{DecisionTreeClassifier} with
$\mathit{max\_depth} = 4$; for Breast, Heart, and Steel we use pre-trained
trees from the SortingHat benchmark~\cite{sortinghat}.

\textbf{Evaluation protocol.}  For each dataset we run three modes:
\emph{Hard clear} (exact comparisons), \emph{Soft adaptive} (per-node
degrees via Eq.~\eqref{eq:degmap}), and \emph{Adaptive HE}.  HE inference is
executed in a containerised environment (Podman, Ubuntu 22.04) on a laptop
CPU.  We report accuracy, the adaptive approximation gap with respect to the
hard baseline, the residual adaptive HE gap with respect to adaptive
cleartext, and average per-sample HE inference time.

\textbf{Implementation setup used in the POC.}  The proof-of-concept
evaluates an already trained tree at runtime rather than retraining a model
online.  In practice, the C++ pipeline loads a tree from \texttt{tree.csv},
\texttt{tree.json}, or compatible benchmark formats, reconstructs the binary
tree in memory, and then runs both cleartext and HE inference.  The full
pipeline is orchestrated by a PowerShell script (\texttt{run.ps1}) that
builds and runs two executables, \texttt{poc\_clear} and \texttt{poc\_he},
inside a Podman container with Microsoft SEAL 4.1.2.  Unless otherwise
stated, we evaluate the first $n$ samples of the test split, with
$n \in [20,32]$ in the main reported runs.  For the two Spam-family
datasets, the current campaign only validates small cleartext sanity checks:
the real-HE pipeline still runs out of memory under the available Podman
configuration.

\textbf{HE runtime profiles.}  The current implementation selects one of
three CKKS profiles from the feature dimension $d$.  For $d \leq 4$, a
\emph{small} profile uses layout resolution $n_{\max}=4$, global polynomial
degree $8$, maximum adaptive degree $32$, and scale modulus size $50$ bits.
For $5 \leq d \leq 16$, a \emph{medium} profile uses $n_{\max}=2$, global
degree $4$, maximum adaptive degree $16$, and scale modulus size $45$ bits.
For $d > 16$, a \emph{large} profile uses $n_{\max}=2$, global degree $4$,
maximum adaptive degree $8$, and scale modulus size $40$ bits.  The target
multiplicative depth is derived automatically from the largest adaptive
degree in the tree, with a minimum depth of 10 levels.

\begin{table}[t]
\centering
\caption{Runtime HE profiles used in the implementation.}
\label{tab:he_profiles}
\begin{tabular}{@{}l c c c c@{}}
\toprule
Profile & Feature range & $n_{\max}$ & global degree & max adaptive degree \\
\midrule
Small  & $d \leq 4$ & 4 & 8 & 32 \\
Medium & $5 \leq d \leq 16$ & 2 & 4 & 16 \\
Large  & $d > 16$ & 2 & 4 & 8 \\
\bottomrule
\end{tabular}
\end{table}

\subsection{Benchmark Scope}

In addition to the main experiments of our implementation, we collected local
benchmark traces from three adjacent baselines: an Akavia-style SEAL
pipeline, the SortingHat benchmark family, and a Concrete-ML prototype.  We
include these runs in the experimental scope of the paper because they help
position our POC relative to existing implementation styles, but they should
not be interpreted as perfectly apples-to-apples comparisons.  In
particular, Akavia and our POC both rely on real homomorphic execution with
Microsoft SEAL, whereas the Concrete-ML runs available in our workspace are
reported in FHE-simulated mode.  Likewise, the SortingHat runs often use
external pre-trained trees and expose different failure modes, including
out-of-memory cases on larger datasets.

\begin{table}[t]
\centering
\caption{Benchmark families considered in the local experimental campaign.}
\label{tab:benchmark_families}
\footnotesize
\begin{tabular}{@{}p{2.2cm} p{2.8cm} p{4.0cm} p{3.0cm}@{}}
\toprule
Family & HE backend / mode & Main datasets & Role in this paper \\
\midrule
My POC & SEAL, real HE & Iris, Cancer, Wine, Digits, Breast, Heart, Steel & Primary quantitative study \\
Akavia-style baseline & SEAL, real HE & Iris, Cancer, Wine, Digits & Historical packed baseline \\
SortingHat & External benchmark pipelines, real HE & Breast, Heart, Steel, Spam-family & Stress-test and OOM analysis \\
Concrete-ML & Concrete-ML, FHE simulation & Iris, Cancer, Wine, Digits, Breast, Heart, Steel & Adjacent reference only \\
\bottomrule
\end{tabular}
\normalsize
\end{table}

The quantitative tables that follow focus primarily on \emph{My POC},
because this is the only family in which we control the full inference
pipeline, including hard, adaptive cleartext, and adaptive HE variants under
a common evaluation protocol.  Akavia results are used as a historical SEAL
baseline on standard datasets; SortingHat results are used to document
behaviour on externally supplied trees and memory-intensive cases; and
Concrete-ML results are used only as an auxiliary point of comparison since
they correspond to FHE-simulated execution rather than the same real-HE
runtime used in the rest of the paper.  For these external families, the
comparative numbers reported below are taken from the locally generated
aggregate benchmark plots available in our workspace and therefore cover
only the dataset-project pairs that are visible in those plots.

\subsection{Main Results}

Table~\ref{tab:main_results} reports the full experimental results on the
representative batch sizes used in our evaluation.

\begin{table}[t]
\centering
\caption{Accuracy (\%) across the main inference modes retained in this
paper.  Timing is average adaptive HE inference time per sample (ms).}
\label{tab:main_results}
\setlength{\tabcolsep}{4pt}
\begin{tabular}{@{}l r r r r r@{}}
\toprule
Dataset & $d$ & $n$ & Hard & Soft Ad. & HE Ad. & ms/samp. \\
\midrule
Iris    & 4  & 20  & 100.0 &  95.0 &  95.0 & 286 \\
Cancer  & 30 & 32  & 100.0 & 100.0 &  96.9 & 151 \\
Wine    & 13 & 26  & 100.0 & 100.0 & 100.0 & 224 \\
Digits  & 64 & 32  & 100.0 &  96.9 &  96.9 & 146 \\
Breast  & 30 & 32  &  90.6 &  90.6 &  90.6 & 320 \\
Heart   & 13 & 26  &  84.6 &  84.6 &  84.6 &  91 \\
Steel   & 33 & 32  & 100.0 & 100.0 & 100.0 &  95 \\
\bottomrule
\end{tabular}
\end{table}

\textbf{Precision decomposition.}  Table~\ref{tab:decomposition} presents
the decomposed accuracy loss, showing the dominant error source for each
dataset.

\begin{table}[t]
\centering
\caption{Decomposed adaptive accuracy loss. $\Delta_{\mathrm{approx}}$ is
the adaptive cleartext gap versus hard; $\Delta_{\mathrm{CKKS}}$ is the
additional adaptive HE gap versus adaptive cleartext. Regime per
Section~\ref{sec:regimes}.}
\label{tab:decomposition}
\begin{tabular}{@{}l c c c c@{}}
\toprule
Dataset & $\mathit{min\_gap}$ & Regime
        & $\Delta_{\mathrm{approx}}$ (\%) & $\Delta_{\mathrm{CKKS}}$ (\%) \\
\midrule
Iris    & 0.017 & II       & 5.0  & 0.0  \\
Cancer  & 0.002 & III      & 0.0  & 3.1 \\
Wine    & 0.003 & III      & 0.0  & 0.0  \\
Digits  & --    & --       & 3.1  & 0.0  \\
Breast  & $<0.02$ & III    & 0.0 & 0.0 \\
Heart   & $<0.03$ & III    & 0.0 & 0.0 \\
Steel   & $<0.02$ & III    & 0.0 & 0.0  \\
\bottomrule
\end{tabular}
\end{table}

\subsection{Cross-Project Benchmark Comparison}

Table~\ref{tab:benchmark_accuracy_results} summarises the accuracy values
visible in the local cross-project benchmark graph.  We use the adaptive HE
result of our POC as the main comparison point because it reflects the final
configuration after node-wise normalisation.  Digits is absent from this
accuracy comparison because no aligned external value is displayed in the
available graph.

\begin{table}[t]
\centering
\caption{Cross-project accuracy comparison (\%).  ``--'' indicates that no
value is displayed for that project-dataset pair in the local benchmark
graph.}
\label{tab:benchmark_accuracy_results}
\footnotesize
\begin{tabular}{@{}l c c c c@{}}
\toprule
Dataset & My POC HE Ad. & Akavia & SortingHat & Concrete-ML \\
\midrule
Cancer & 96.9 & 78.1 & --   & 93.3 \\
Wine   & 100.0 & 61.5 & --   & 100.0 \\
Iris   & 95.0 & 90.0 & --   & 100.0 \\
Breast & 90.6 & --   & 90.6 & 100.0 \\
Heart  & 84.6 & --   & 81.8 & 81.8 \\
Steel  & 100.0 & --   & 100.0 & 100.0 \\
\bottomrule
\end{tabular}
\normalsize
\end{table}

Table~\ref{tab:benchmark_latency_results} reports the average per-sample
times extracted from the local latency benchmark graph.  Because some
projects were run with different sample counts, we report the visible batch
size $n$ alongside each timing when it is legible in the graph.

\begin{table}[t]
\centering
\caption{Cross-project average HE inference time per sample (ms), with the
visible benchmark batch size in parentheses.}
\label{tab:benchmark_latency_results}
\footnotesize
\begin{tabular}{@{}p{1.5cm} p{2.2cm} p{2.2cm} p{2.2cm} p{2.2cm}@{}}
\toprule
Dataset & My POC & Akavia & SortingHat & Concrete-ML \\
\midrule
Breast & 320.3 (32) & -- & 154.6 & 190.5 \\
Cancer & 151.1 (32) & 18996.3 & -- & 430.5 \\
Heart  & 90.5 (26) & -- & 35.1 & 537.1 \\
Iris   & 286.1 (20) & 12695.2 (10) & -- & 561.3 (10) \\
Steel  & 95.4 (32) & -- & 38.8 & 418.2 \\
Wine   & 224.5 (26) & 13274.6 & -- & 417.7 \\
\bottomrule
\end{tabular}
\normalsize
\end{table}

\textbf{Adaptive pipeline behaviour.}  The adaptive configuration is now the
central reference point of the paper.  On Breast, Heart, and Steel,
node-wise logit normalisation allows adaptive cleartext and adaptive HE to
match the hard baseline exactly.  On Wine, the adaptive pipeline also
preserves the hard-tree accuracy under HE.  On Digits, adaptive HE matches
adaptive cleartext at 96.9\%.  Cancer remains the main counter-example:
adaptive cleartext reaches 100.0\%, but adaptive HE still drops to 96.9\%,
showing that very small gaps can remain CKKS-sensitive even after local
normalisation.

\subsection{Precision-Regime Observations}

\textbf{Iris (Regime II).}  The minimum gap $g = 0.017$ lies well above
$\varepsilon = 10^{-6}$ and slightly below $\delta = 0.05$.  In the latest
run, the adaptive pipeline stays 5 points below the hard baseline
(95.0\% vs.\ 100.0\%) both in cleartext and in HE.  This suggests that, on
this small batch, the current degree heuristic is slightly too aggressive.

\textbf{Cancer (Regime III).}  The minimum gap $g \approx 0.002$ is far
below $\delta$.  After node-wise normalisation, the adaptive cleartext run
now reaches 100.0\%, so the approximation loss disappears on the evaluated
batch.  The remaining degradation is HE-specific: adaptive HE drops to
96.9\%, yielding $\Delta_{\mathrm{CKKS}} \approx 3.1\%$.  Cancer is
therefore the clearest current example of a dataset that remains limited by
very small gaps even after the soft-comparison calibration has been fixed.

\textbf{Wine (Regime III).}  Despite $g = 0.003 < \delta$, cleartext soft
adaptive inference achieves 100\% accuracy, and adaptive HE preserves that
result.  Wine therefore remains a strong example of adaptive degree
selection improving robustness without any visible cleartext loss.

\textbf{Heart, Breast, and Steel after node-wise logit normalisation.}
For the three pre-trained SortingHat trees, the dominant source of loss was
not CKKS noise itself but the lack of local logit normalisation in the soft
comparison stage.  After calibrating each internal node with a node-specific
normalisation factor $\max_x |x[\mathrm{feat}] - \theta|$ and assigning the
polynomial degree from the normalised gap $\mathit{min\_gap}/\mathit{norm}$,
the adaptive soft mode recovers the hard-tree accuracy in cleartext and under
CKKS on all three datasets.  Concretely, Breast rises from 40.6\% to 90.6\%,
Heart from 57.7\% to 84.6\%, and Steel from 62.5\% to 100.0\% in adaptive
mode, with HE matching the cleartext adaptive baseline in each case.  This
shows that these datasets were not inherently CKKS-limited; instead, they
were poorly conditioned in the unnormalised comparison space.

\subsection{Runtime Performance}

HE inference time ranges from 90 ms/sample (Heart, $d = 13$, batch 26) to
320 ms/sample (Breast, $d = 30$, batch 32) in the adaptive mode among the
latest validated runs, with Iris still remaining costlier than the larger
datasets because of key-switching overhead at small batch sizes.  These
timings are consistent with previously reported
results for CKKS-based ML inference on laptop hardware~\cite{dathathri2019chet}
and confirm that private decision tree inference is practical for
non-real-time applications.

\subsection{Cross-Benchmark Observations}

The additional benchmark families available in our workspace support three
main observations.  First, the Akavia-style SEAL baseline is substantially
slower than our current implementation on the overlapping standard datasets:
for example, the local graph reports 12.7~s/sample on Iris, 18.9~s/sample on
Cancer, and 13.3~s/sample on Wine, versus 286~ms, 151~ms, and 224~ms/sample
respectively for our POC.  Accuracy is also lower on the same overlaps, with
78.1\% on Cancer, 61.5\% on Wine, and 90.0\% on Iris compared with our
adaptive HE results of 96.9\%, 100.0\%, and 95.0\%.

Second, the SortingHat benchmark family is competitive on the externally
supplied trees from which it originates.  In the local comparison graph it
matches our adaptive HE accuracy on Breast (90.6\%) and Steel (100.0\%), and
trails slightly on Heart (81.8\% vs.\ 84.6\%).  It is also faster on those
tasks, with visible per-sample times of 154.6~ms on Breast, 35.1~ms on
Heart, and 38.8~ms on Steel, although these comparisons remain imperfect
because the benchmarked sample counts are not always identical.

Third, the Concrete-ML prototype provides a strong neighbouring reference in
accuracy while remaining a simulation-based backend.  In the local graph it
reaches 100.0\% on Wine, Breast, and Steel, 93.3\% on Cancer, and 81.8\% on
Heart.  Its latency remains above our POC on the displayed overlaps, ranging
from about 190~ms/sample on Breast to about 561~ms/sample on Iris.  These
results are useful for situating our implementation, but they should not be
be read as a like-for-like real-HE comparison.

%% =========================================================
\section{Related Work}
\label{sec:related}

\textbf{Private tree inference.}  Bos \emph{et al.}~\cite{bos2014private}
pioneered FHE-based decision tree evaluation using BGV.  Tai \emph{et
al.}~\cite{tai2017mining} proposed a garbled-circuit approach with lower
communication.  Kiss \emph{et al.}~\cite{kiss2019sok} provided a
comparative survey.  Tueno \emph{et al.}~\cite{tueno2019private} exploited
TFHE for exact binary evaluation, avoiding approximation entirely but at the
cost of much higher per-gate overhead.

\textbf{Polynomial approximation for FHE.}  Han and Ki~\cite{han2020better}
formalised BSGS polynomial evaluation for CKKS.  Chiang \emph{et
al.}~\cite{chiang2022} studied Chebyshev-based approximations of the
comparison function.  Compared to these works, we evaluate approximations at
a per-node granularity and explicitly quantify the approximation contribution
to downstream task accuracy.

\textbf{HBDT-SumPath.}  Shin \emph{et al.}~\cite{shin2022privacy} introduced
the SumPath depth-reduction algorithm that we build upon.  Our work extends
it with an adaptive degree selection strategy and a formal precision
decomposition framework absent in the original paper.

\textbf{Precision analysis of CKKS.}  Cheon \emph{et
al.}~\cite{cheon2021efficient} and Bossuat \emph{et al.}~\cite{bossuat2021}
studied error propagation in CKKS for polynomial evaluation circuits.  Our
analysis is application-specific: we connect CKKS noise to ML accuracy at
the semantic level, through the lens of min\_gap, a previously unreported
metric for privacy-preserving decision tree evaluation.

%% =========================================================
\section{Conclusion}
\label{sec:conclusion}

We have presented a complete implementation of private decision tree
inference under CKKS and, more importantly, a principled decomposition of
the accuracy loss into two independent sources: polynomial approximation
error and CKKS residual noise.  The \emph{minimum feature gap}
$\mathit{min\_gap}$ emerges as the key predictor of which source dominates:
datasets outside the most critical small-gap regime show little or no
additional CKKS loss, while very small-gap datasets remain sensitive to
calibration and modelling choices.  In particular, our experiments show
that part of the loss that
initially appeared to be HE-specific was in fact caused by unnormalised
comparison logits.  After introducing node-wise logit normalisation, the
adaptive pipeline fully recovers the hard-tree baseline on Breast, Heart,
and Steel in both cleartext and CKKS, while Cancer remains the main dataset
whose adaptive HE accuracy still trails the adaptive cleartext baseline.

Our results have two immediate practical implications.  First,
practitioners deploying private inference should compute $\mathit{min\_gap}$
before selecting CKKS parameters: a small gap demands larger scale
$\Delta = 2^{50}$ or additional multiplicative depth only after the
comparison space has been properly calibrated.  Second, model trainers should
consider both $\mathit{min\_gap}$ and node-wise logit scale as regularisation
targets alongside standard accuracy metrics, so that the resulting tree is
cryptographically friendly before encryption.

Future work will investigate min\_gap-regularised training, extension to
random forests and gradient-boosted trees, stronger calibration strategies
for externally supplied trees, dedicated low-memory HE profiles for the
Spam-family datasets, and a formal worst-case bound relating
$\mathit{min\_gap}$, node-wise normalisation, CKKS parameters, and
guaranteed accuracy.

\paragraph{Acknowledgements.}
The authors thank the OpenFHE and Microsoft SEAL communities for their
open-source implementations.

%% =========================================================
\begin{thebibliography}{16}

\bibitem{albrecht2021homomorphic}
Albrecht, M., et al.:
Homomorphic encryption standard.
In: Lauter, K. et al. (eds.) Protecting Privacy through Homomorphic Encryption, pp.~31--62.
Springer, Cham (2021).
\href{https://doi.org/10.1007/978-3-030-77287-1_2}
     {https://doi.org/10.1007/978-3-030-77287-1\_2}

\bibitem{bos2014private}
Bos, J.W., Lauter, K., Loftus, J., Naehrig, M.:
Improved security for a ring-based fully homomorphic encryption scheme.
In: Stam, M. (ed.) IMACC 2013, LNCS, vol.~8308, pp.~45--64.
Springer, Heidelberg (2013).
\href{https://doi.org/10.1007/978-3-642-45239-0_4}
     {https://doi.org/10.1007/978-3-642-45239-0\_4}

\bibitem{bossuat2021}
Bossuat, J.P., Mouchet, C., Troncoso-Pastoriza, J., Hubaux, J.P.:
Efficient bootstrapping for approximate homomorphic encryption with
non-sparse keys.
In: Canteaut, A., Standaert, F.X. (eds.) EUROCRYPT 2021, LNCS, vol.~12696,
pp.~587--617. Springer, Cham (2021).

\bibitem{chen2020simple}
Chen, H., Iliashenko, I., Laine, K.:
When HEAAN meets SEAL: High-performance evaluation of high-precision rational polynomials.
Cryptology ePrint Archive 2019/1325 (2019).

\bibitem{cheon2017homomorphic}
Cheon, J.H., Kim, A., Kim, M., Song, Y.:
Homomorphic encryption for arithmetic of approximate numbers.
In: Takagi, T., Peyrin, T. (eds.) ASIACRYPT 2017, LNCS, vol.~10624,
pp.~409--437. Springer, Cham (2017).
\href{https://doi.org/10.1007/978-3-319-70694-8_15}
     {https://doi.org/10.1007/978-3-319-70694-8\_15}

\bibitem{cheon2021efficient}
Cheon, J.H., Han, K., Kim, A., Kim, M., Song, Y.:
Bootstrapping for approximate homomorphic encryption.
In: Nielsen, J.B., Rijmen, V. (eds.) EUROCRYPT 2018, LNCS, vol.~10820,
pp.~360--384. Springer, Cham (2018).

\bibitem{chiang2022}
Chiang, Y.C., et al.:
Comparison and decision-tree homomorphic evaluation with CKKS.
Cryptology ePrint Archive 2022/280 (2022).

\bibitem{dathathri2019chet}
Dathathri, R., et al.:
CHET: An optimizing compiler for fully-homomorphic neural-network inferencing.
In: Proceedings of PLDI 2019, pp.~142--156. ACM, New York (2019).

\bibitem{han2020better}
Han, K., Ki, D.:
Better bootstrapping for approximate homomorphic encryption.
In: Jarecki, S. (ed.) CT-RSA 2020, LNCS, vol.~12006, pp.~364--390.
Springer, Cham (2020).

\bibitem{kiss2019sok}
Kiss, \'{A}., Liu, J., Schneider, T., Asokan, N., Pinkas, B.:
Private set intersection for unequal set sizes with mobile applications.
PoPETs 2017(4), 177--197 (2017).

\bibitem{shin2022privacy}
Shin, Y., Choi, R., Cheon, J.H.:
Privacy-preserving decision tree inference with CKKS.
Cryptology ePrint Archive 2021/1337 (2021).

\bibitem{sortinghat}
Oliynyk, D., Mayer, R., Rauber, A.:
I know what you trained last summer: A survey on stealing machine learning models and defences.
ACM Comput. Surv. 55(14), 1--41 (2023).

\bibitem{tai2017mining}
Tai, R.K.H., Ma, J.P.K., Zhao, Y., Chow, S.S.M.:
Privacy-preserving decision trees evaluation via linear functions.
In: Foley, S.N., Gollmann, D., Snekkenes, E. (eds.) ESORICS 2017, LNCS,
vol.~10493, pp.~494--512. Springer, Cham (2017).

\bibitem{tueno2019private}
Tueno, A., Kerschbaum, F., Katzenbeisser, S.:
Private evaluation of decision trees using sublinear cost.
PoPETs 2019(1), 266--286 (2019).

\end{thebibliography}

\end{document}
